\section{Fundamental Concepts}

\paragraph{Supervised Learning}
Supervised learning deals with sets of iid observations $(x_{i}, y_{i})$. The first element is called the data and the second is called the label, and these live in data space $\mathcal{X}$ and label space $\mathcal{Y}$ respectively. The data and labels are sampled from a probability distribution $P$ defined on $\mathcal{X}\times\mathcal{Y}$.

Furthermore, the labels are generated from the data by some unknown function $f$, the estimation of which is the goal of supervised learning. The estimate is chosen from some parametrized class $\mathcal{F}$, which consists of members $f_{\theta}$. The performance of each choice is evaluated based on expected performance on samples from $P$ as measured by a loss function $L$. That is, the performance $\mathcal{R}(\tilde{f})$ is given by
\begin{align}
  \mathcal{R}(\tilde{f}) = \expval{L(\tilde{f}(X), f(X))}_{P}.
\end{align}

\paragraph{The Curse of Dimensionality}
The curse of dimensionality refers to the set of phenomena where obtaining an adequateand regular fit to a dataset generally requires a number of samples that increases exponentially with the data dimensionality. Given that many machine learning applications concern high-dimensional data, this is a serious obstacle to doing machine learning.

\paragraph{Breaking the Curse}
Neural networks can break the curse of dimensionality by imposing sparsity on the weights, thus providing regularity. This comes, however, at the cost of making assumptions about the target function, which is undesirable. Geometric deep learning can instead provide a different source of regularity by analyzing structure and geometry in the data.

\paragraph{The Geometry of Data}
To start analysing data in terms of geometry, we will decouple the signal domain as follows: We take it that the data domain $\mathcal{X}$ is composed of some domain $\Omega$ and some signal which maps points in $\Omega$ to some other domain $\mathcal{C}$. That is,
\begin{align}
  \mathcal{X} = \mathcal{X}(\Omega, \mathcal{C}) = \left\{x: \Omega\to\mathcal{C}\right\}.
\end{align}

The canon example for this is an $m\times n$-size RGB image. In this cse we would have $\Omega = \mathbb{Z}_{m}\times\mathbb{Z}_{n}$ and $\mathcal{C} = \mathbb{R}^{3}$ (or some other three-dimensional domain). In that case, a particular piece of data (a specific image) is a particular signal which describes how you need to distribute your RGB values to form a particular image.

The purpose of doing this is to separate out the geometry of the data. When considering for instance mirror or translation symmetry in an image, that only applies to the spatial part of the image, not the color part. Another thing it achieves is to separate out a part with the structure of a vector space. For instance, we can define an inner product between the RGB values of two different pixels to determine whether they are similar. Adding a measure $\mu$ on $\Omega$, we can extend that to define an inner product between signals according to
\begin{align}
  \braket{x_{1}}{x_{2}} = \inte{\Omega}{}\dd{\mu(\omega)}\braket{x_{1}(\omega)}{x_{2}(\omega)}.
\end{align}

\paragraph{Geometric Priors}
There are two principles which we will refer to as our geometric priors: symmetry and scale separation.

Symmetry will here refer to the physical notion of symmetry, and thus include both discrete and continuous symmetries. Separating out the geometric part $\Omega$ of the data domain allows us to work directly with its symmetries.

Scale separation is the notion that phenomena at different scales can be treated separately. For instance, returning to the image analogy, scale separation entails that instead of always treating the minutia of a sample image, we can apply coarse-graining operations such as pooling without losing essential information.
